% ============================================================================
%  tmux-cheatsheet.tex
%  Quick reference for tmux key bindings and commands (default bindings).
%  Requires linuxnotes.sty in the same folder.
%  Compile twice with pdflatex so "Page X of Y" resolves.
% ============================================================================
\documentclass[11pt]{article}

\usepackage[pagenumbers]{linuxnotes}

\linuxnotesmeta{tmux Cheat Sheet}

% Header row style shared by every table.
\newcommand{\headrow}{\rowcolor{navy!8}}

% \key from linuxnotes.sty cannot take a percent sign, so \keylit is the same
% boxed style with its argument typeset as written (use \% for a percent).
\newcommand{\keylit}[1]{%
  \fcolorbox{rulecol}{keybg}{\rule[-0.2em]{0pt}{0.9em}\footnotesize\texttt{#1}}%
}
% \tk{d} shows the prefix followed by a key: Ctrl+b, then d.
\newcommand{\tk}[1]{\key{Ctrl+b}\enspace\keylit{#1}}

\begin{document}

\linuxnotestitle{WORK-DEV}{tmux Cheat Sheet}{Sessions, windows, panes, and scrolling with the default key bindings}

\section{The prefix key}

Every shortcut starts with the prefix \key{Ctrl+b}. In the tables, \tk{d} means press
\key{Ctrl+b}, release both keys, then press \key{d}. If nothing happens, you probably
held \key{Ctrl} down too long.

\section{Sessions}

A session is a whole workspace. It keeps running after you detach until you kill it.

\par\noindent\begin{minipage}{\linewidth}\centering
\captionof{table}{Session keys and commands.}
\renewcommand{\arraystretch}{1.2}
\begin{tabular}{>{\raggedright\arraybackslash}p{0.32\linewidth} >{\raggedright\arraybackslash}p{0.27\linewidth} >{\raggedright\arraybackslash}p{0.31\linewidth}}
\toprule
\headrow
\bfseries Action & \bfseries Inside tmux & \bfseries From a terminal \\
\midrule
Detach and leave it running & \tk{d} & None \\
List sessions & \tk{s} & \cmd{tmux ls} \\
Attach to a session & None & \cmd{tmux attach -t} \var{NAME} \\
Start a new session & None & \cmd{tmux new -s} \var{NAME} \\
Rename the current session & \tk{\$} & None \\
Close the current session & \tk{:} \cmd{kill-session} & \cmd{tmux kill-session -t} \var{NAME} \\
Close every session & None & \cmd{tmux kill-server} \\
\bottomrule
\end{tabular}
\end{minipage}\par\medskip

\section{Windows}

A window is like a browser tab; the status bar lists windows by number.

\par\noindent\begin{minipage}{\linewidth}\centering
\captionof{table}{Window keys.}
\renewcommand{\arraystretch}{1.2}
\begin{tabular}{>{\raggedright\arraybackslash}p{0.55\linewidth} >{\raggedright\arraybackslash}p{0.35\linewidth}}
\toprule
\headrow
\bfseries Action & \bfseries Keys \\
\midrule
New window & \tk{c} \\
Next window & \tk{n} \\
Previous window & \tk{p} \\
Jump to a window by number & \tk{0} through \keylit{9} \\
Pick from a list of all windows & \tk{w} \\
Toggle to the last window used & \tk{l} \\
Rename the current window & \tk{,} \\
Close the current window (asks first) & \tk{\&} \\
\bottomrule
\end{tabular}
\end{minipage}\par\medskip

\section{Panes}

Panes split one window into several terminals. \tk{z} zooms one pane to full size and back.

\par\noindent\begin{minipage}{\linewidth}\centering
\captionof{table}{Pane keys.}
\renewcommand{\arraystretch}{1.2}
\begin{tabular}{>{\raggedright\arraybackslash}p{0.46\linewidth} >{\raggedright\arraybackslash}p{0.44\linewidth}}
\toprule
\headrow
\bfseries Action & \bfseries Keys \\
\midrule
Split left and right & \tk{\%} \\
Split top and bottom & \tk{"} \\
Move to the pane in a direction & \key{Ctrl+b} then an arrow key \\
Cycle through panes & \tk{o} \\
Show pane numbers, then press one to jump & \tk{q} \\
Zoom the current pane in or out & \tk{z} \\
Resize the current pane & \key{Ctrl+b}, then \key{Ctrl} plus an arrow key \\
Cycle through preset layouts & \tk{Space} \\
Turn the current pane into its own window & \tk{!} \\
Close the current pane (asks first) & \tk{x} \\
\bottomrule
\end{tabular}
\end{minipage}\par\medskip

\section{Scrolling, copying, and commands}

By default the mouse wheel does not scroll, so use copy mode; press \key{q} to leave it.

\par\noindent\begin{minipage}{\linewidth}\centering
\captionof{table}{Copy mode and command keys.}
\renewcommand{\arraystretch}{1.2}
\begin{tabular}{>{\raggedright\arraybackslash}p{0.46\linewidth} >{\raggedright\arraybackslash}p{0.44\linewidth}}
\toprule
\headrow
\bfseries Action & \bfseries Keys \\
\midrule
Enter copy mode (scroll back) & \tk{[} \\
Scroll in copy mode & Arrow keys, \key{PgUp}, \key{PgDn} \\
Search backward in copy mode & \key{?} then the text, then \key{Enter} \\
Start selecting text & \key{Space} \\
Copy the selection and leave copy mode & \key{Enter} \\
Paste the copied text & \tk{]} \\
Open the tmux command prompt & \tk{:} \\
Show every key binding & \tk{?} (press \key{q} to close) \\
\bottomrule
\end{tabular}
\end{minipage}\par\medskip

\notebox{Tip.}{For mouse scrolling and clicking between panes, add \cmd{set -g mouse on}
to \file{~/.tmux.conf}, then run \cmd{tmux source-file ~/.tmux.conf}.}

\vspace{1em}
\noindent{\footnotesize\color{navysoft}\copyright\ 2026 Jason Powell. Part of Work-Dev. Licensed under
\href{https://creativecommons.org/licenses/by/4.0/}{Creative Commons Attribution 4.0 (CC BY 4.0)}.}

\end{document}
